\section{Theorie}
\label{sec:Theorie}

\subsection{Das ideale Gasmodell}
\label{subsec:IdealesGasmodell}

    In dem idealen Gasmodell wird ein Gas als Ensemble von nicht wechselwirkenden, punktförmigen Teilchen betrachtet. 
    Diese Teilchen bewegen sich zufällig und stoßen elastisch miteinander sowie mit den Wänden des Behälters, in dem sich das Gas befindet. 
    Wenn das System im thermischen Gleichgewicht ist, folgt die Geschwindigkeit der Teilchen der Boltzmann-Verteilung.
    Daraus ergibts sich die Zustandsgleichung des idealen Gases:
    \begin{equation}
        pV = \nu RT,
        \label{eq:idealesGas}
    \end{equation}
    wobei \(p\) der Druck, \(V\) das Volumen, \(\nu\) die Stoffmenge in Mol, \(R\) die universelle Gaskonstante und \(T\) die absolute Temperatur ist.
    Das ideale Gasmodell ist eine gute Näherung für reale Gase bei niedrigen Drücken und hohen Temperaturen, 
    wo die Wechselwirkungen zwischen den Molekülen vernachlässigbar sind. 

\subsection{Die Van-der-Waals-Gleichung}
\label{subsec:VanDerWaalsGleichung}

    Die Van-der-Waals-Gleichung ist eine Modifikation der idealen Gasgleichung, die die endliche Größe der Moleküle und die Anziehungskräfte zwischen ihnen berücksichtigt. 
    Sie lautet:
    \begin{equation}
        \left( p + a \frac{\nu^2}{V^2} \right) (V - \nu b) = nRT,
        \label{eq:VanDerWaals}
    \end{equation}
    wobei \(a\) und \(b\) empirische Konstanten sind, die die Wechselwirkungen zwischen den Molekülen und das Eigenvolumen der Moleküle repräsentieren.
    Der Term \(a \frac{n^2}{V^2}\) korrigiert den Druck um die Anziehungskräfte, während der Term \(\nu b\) das Volumen um das Eigenvolumen der Moleküle reduziert.
    Diese Gleichung bietet eine genauere Beschreibung des Verhaltens realer Gase. 
    